\section{The growing-band velocity profile}
\label{sec:terminal-modal-velocity-profile}

The velocity target has a simpler source kernel than the position target.
For a source $F$ injected at time $n+1$, put $\ell=m-n$.  Subtracting the two
Green responses at times $m+1$ and $m$ gives the exact identity
\begin{equation}
 \begin{split}
 &r^\ell\frac{\sin((\ell+1)\phi)}{\sin\phi}F
 -(1-q)r^{\ell-1}\frac{\sin(\ell\phi)}{\sin\phi}F\\
 &\hspace{35mm}=(1-q)r^{\ell-1}\cos((\ell+1)\phi)F.
 \end{split}
 \label{eq:terminal-modal-velocity-source-kernel}
\end{equation}
Thus no small-sine factor survives.  We use the position-limit functions from
\eqref{eq:terminal-modal-position-limit-profile}.  Put
\begin{equation}
 p_\star=p_0(1),\qquad
 \mu(x)=\frac{p_0(x)}2\frac{1-t_x^2}{1+t_x^2}+\frac25,
 \label{eq:terminal-modal-velocity-limit-data}
\end{equation}
and
\begin{equation}
 h(x)=e^{-(1-x)/16}u(x),qquad
 g(x)=e^{-(1-x)/16}\mu(x).
 \label{eq:terminal-modal-velocity-amplitudes}
\end{equation}

\begin{lemma}[Uniform velocity limit; proved here]
\label{lem:terminal-modal-velocity-limit}
Uniformly for $1\le s\le H_m$,
\begin{equation}
 \frac m\alpha\left(V_{2s}+\frac{G_{2s}}5\right)
 =T_s+o(1),
 \label{eq:terminal-modal-velocity-uniform-limit}
\end{equation}
where
\begin{align}
 T_s={}&-\frac{p_\star}{6}+\frac{2}{15}
 +8(-1)^s\int_0^1h'(x)
 \left\{\cos(\pi sx)+\frac13\cos(3\pi sx)\right\}\,dx
 \notag\\
 &+\frac{(-1)^s}{16}\int_0^1g(x)
 \left\{\cos(\pi sx)+\cos(3\pi sx)\right\}\,dx.
 \label{eq:terminal-modal-velocity-limit-functional}
\end{align}
\end{lemma}

\begin{proof}
We evaluate the exact five pieces following
\cref{lem:terminal-modal-source-scalar-reduction}.  First consider all
derivative sources $U_n$ together.  Equations
\eqref{eq:terminal-modal-position-U-limit},
\eqref{eq:terminal-modal-position-source-Taylor}, and
\eqref{eq:terminal-modal-velocity-source-kernel} give
\begin{equation}
 \frac{R^U_{m,s}-(1-q)R^{U,-}_{m,s}}{q^3}
 =16\pi s(-1)^s\int_0^1h(x)
 \sin(2\pi sx)\cos(\pi sx)\,dx+o(1).
 \label{eq:terminal-modal-velocity-U-limit}
\end{equation}
Here $R^{U,-}_{m,s}$ is the same source response one time earlier.  The
error is uniform: the source Taylor error costs $O(\phi_s)$, damping costs
$O(q+m\phi_s^2)$, and the Riemann error costs $O(s/m)$, all $o(1)$ on
$1\le s\le H_m$.

The mass source has transform $2M_n\cos(2\pi s n/m)$.  The exact kernel
\eqref{eq:terminal-modal-velocity-source-kernel}, the uniform convergence
$M_n/q^4=\mu(n/m)+O(q)$ obtained from
\eqref{eq:terminal-modal-finite-mu-M}, and $2mq=1/8$ give
\begin{equation}
 \frac{R^M_{m,s}-(1-q)R^{M,-}_{m,s}}{q^3}
 =\frac{(-1)^s}{8}\int_0^1g(x)
 \cos(2\pi sx)\cos(\pi sx)\,dx+o(1).
 \label{eq:terminal-modal-velocity-M-limit}
\end{equation}
The same error estimates are uniform on the growing band.

We next compute the two boundary pieces.  At $q=0$, the first two correction
vectors in the fixed early replay give
\begin{equation}
 c_3=\left(\frac7{20},\frac1{10},\frac9{40}\right),
 \qquad
 c_4=\left(\frac14,\frac{17}{80},\frac18,\frac{27}{80}\right).
 \label{eq:terminal-modal-velocity-early-corrections}
\end{equation}
Their even-transform masses are $1$ and $8/5$.  Since
$\phi_s^2=o(q)$ uniformly on the band, the corresponding finite-$q$
transforms are $q^3(1+o(1))$ and $q^3(8/5+o(1))$.  Applying
\eqref{eq:terminal-modal-velocity-source-kernel} to the homogeneous base and
including its assigned $q mG_{2s}/5$ term gives
\begin{equation}
 \frac{B^V_{m,s}}{q^3}
 =(-1)^{s+1}\frac45e^{-1/16}+o(1).
 \label{eq:terminal-modal-velocity-base-limit}
\end{equation}
Indeed, the correction part contributes
$(-1)^{s+1}(3/5)e^{-1/16}$, the ideal homogeneous part cancels exactly in
$A_{m+1}-(1-q)A_m$, and the assigned ideal term contributes the remaining
$(-1)^{s+1}e^{-1/16}/5$.

For the final boundary, put $S_n=\sigma_n/q^3$.  The exact frontier recurrence
in \eqref{eq:terminal-modal-frontier-average-residual} is
\begin{equation}
 S_n=\frac{P_n/q^2+\eta S_{n-1}
 -\eta P_{n-1}/(2q^2)-1/5}{1+q}.
 \label{eq:terminal-modal-velocity-frontier-scaled}
\end{equation}
It is a contraction with limiting factor $1/2$.  Since
$P_{m-j}/q^2\to p_\star$ for every fixed $j$, its geometric tail gives
\begin{equation}
 S_m,S_{m-1}\longrightarrow\frac32p_\star-\frac25.
 \label{eq:terminal-modal-velocity-frontier-limit}
\end{equation}
Use the endpoint cancellation
\eqref{eq:terminal-modal-position-endpoint-cancellation}, the final transform
$2(U_m+2U_m+W_m^{\rm end})+o(q^3)$, and the endpoint terms assigned after the
five-piece split.  Substitution of
\eqref{eq:terminal-modal-velocity-frontier-limit} and
$u(1)=(p_\star-4/5)/16$ then gives
\begin{equation}
 \frac{E^V_{m,s}}{q^3}=\frac{p_\star}{2}-\frac25+o(1),
 \label{eq:terminal-modal-velocity-endpoint-limit}
\end{equation}
uniformly in $s$.

It remains only to expose the cancellation among
\eqref{eq:terminal-modal-velocity-U-limit},
\eqref{eq:terminal-modal-velocity-base-limit}, and
\eqref{eq:terminal-modal-velocity-endpoint-limit}.  Product-to-sum and
integration by parts give
\begin{align}
 &16\pi s(-1)^s\int_0^1h(x)\sin(2\pi sx)\cos(\pi sx)\,dx
 \notag\\
 &\quad=\frac{32}{3}\left\{(-1)^sh(0)-h(1)\right\}
 +8(-1)^s\int_0^1h'(x)
 \left\{\cos(\pi sx)+\frac13\cos(3\pi sx)\right\}\,dx.
 \label{eq:terminal-modal-velocity-U-parts}
\end{align}
Because $h(0)=3e^{-1/16}/40$, its parity term cancels
\eqref{eq:terminal-modal-velocity-base-limit}.  Since
$h(1)=(p_\star-4/5)/16$, the remaining boundary term plus
\eqref{eq:terminal-modal-velocity-endpoint-limit} is
$-p_\star/6+2/15$.  Finally use
$\cos(2y)\cos y=(\cos y+\cos3y)/2$ in
\eqref{eq:terminal-modal-velocity-M-limit}.  This proves
\eqref{eq:terminal-modal-velocity-limit-functional} and the uniform limit.
\end{proof}

\begin{theorem}[Growing-band entry velocity profile; proved here]
\label{lem:terminal-modal-velocity-profile}
For all sufficiently large $m$, the literal full-face entry coefficients
satisfy
\begin{equation}
 \frac m\alpha\left|V_{2s}+\frac{G_{2s}}5\right|\le\frac14,
 \qquad 1\le s\le H_m.
 \label{eq:terminal-modal-entry-velocity-profile}
\end{equation}
\end{theorem}

\begin{proof}
We bound the two oscillatory remainders in
\eqref{eq:terminal-modal-velocity-limit-functional} by variation.  Write
$t=e^{-x/8}$.  Since
\begin{equation}
 p_0(x)=\frac25\frac{t^2+10t-1}{1+t^2},
 \label{eq:terminal-modal-velocity-p-rational}
\end{equation}
direct differentiation gives
\begin{equation}
 \left|\frac{dh}{dt}\right|
 =e^{-1/16}\frac{|A(t)|}{80t^{3/2}(1+t^2)^2},qquad
 A(t)=t^4-30t^3+12t^2+10t+3.
 \label{eq:terminal-modal-velocity-h-derivative}
\end{equation}
Now $A''(t)=12t^2-180t+24<0$ on $[7/8,1]$, and
$A'(7/8)<0$.  Hence $A$ is decreasing there, from
$0<A(7/8)=5841/4096<3/2$ to $A(1)=-4$, so $|A(t)|\le4$.  Moreover
$1-e^{-1/8}<1/8$, $t^{3/2}>(7/8)(14/15)=49/60$, and
$(1+t^2)^2>(113/64)^2$.  Integrating
\eqref{eq:terminal-modal-velocity-h-derivative} therefore gives the wholly
rational estimate
\begin{equation}
 \operatorname{TV}(h)
 <\frac18\frac4{80(49/60)(113/64)^2}
 =\frac{1536}{625681}<\frac1{400}.
 \label{eq:terminal-modal-velocity-h-TV}
\end{equation}

For the mass amplitude, direct differentiation shows that the sign of
$g'(x)$ is the sign of
\begin{equation}
 1-10t-9t^2+120t^3+23t^4-30t^5+t^6.
 \label{eq:terminal-modal-velocity-g-sign}
\end{equation}
Its Bernstein coefficients after $t=(7+y)/8$ are
\begin{equation}
 \frac{16854209}{262144},\quad
 \frac{565377}{8192},\quad
 \frac{94677}{1280},\quad
 \frac{202609}{2560},\quad
 \frac{20291}{240},\quad
 \frac{541}{6},\quad96,
 \label{eq:terminal-modal-velocity-g-Bernstein}
\end{equation}
so $g$ is increasing.  Since $g(0)>0$, the mass bound
\eqref{eq:terminal-modal-finite-mass-bound} passes to the limit and gives
\begin{equation}
 \operatorname{TV}(g)<g(1)<\frac{43}{80}.
 \label{eq:terminal-modal-velocity-g-TV}
\end{equation}

Now $0<p_\star<2$, as is immediate from
\eqref{eq:terminal-modal-velocity-p-rational}, so
$|-p_\star/6+2/15|<1/5$.  The $h'$ term in
\eqref{eq:terminal-modal-velocity-limit-functional} is at most
$32\operatorname{TV}(h)/3$.  Integrating each $g$ cosine once by parts makes
its term at most $\operatorname{TV}(g)/(12\pi s)<43/2880$, using $\pi>3$.
Therefore, uniformly for $s\ge1$,
\begin{equation}
 |T_s|<\frac15+\frac{32}{3}\frac1{400}
 +\frac{43}{2880}
 =\frac{3479}{14400}<\frac14.
 \label{eq:terminal-modal-velocity-final-ledger}
\end{equation}
The strict slack is $121/14400$.  The uniform $o(1)$ in
\eqref{eq:terminal-modal-velocity-uniform-limit} proves the theorem for all
sufficiently large $m$.
\end{proof}
